\section{Summation}
In this section, we use the quantum dilogarithm $T_N(z)$ to express $J_N(E;e^{\xi/N})$ without the product of a sequence.
\par
Since the figure-eight knot $E$ is amphicheiral, its colored Jones polynomial is symmetric, that is, $J_{N}(E;q)=J_{N}(E;q^{-1})$.
So we have $J_{N}(E;e^{\xi/N})=J_{N}(E;e^{-\xi/N})$ and we do not need to consider the case where $\xi<0$.
\par
Recall that we choose a positive real number $\theta$ so that $\tan\theta<\frac{\pi}{\xi}$ and that we put $\gamma:=\frac{\xi}{2N\pi\i}$.
Putting $z:=\frac{\xi}{2\pi\i}(1-l/N)$ ($0<l<N$) in \eqref{eq:lem}, we have
\begin{equation*}
  \frac{\exp\left(T_{N}\left(\frac{\xi}{2\pi\i}(1-l/N)-\frac{\xi}{4N\pi\i}\right)\right)}
       {\exp\left(T_{N}\left(\frac{\xi}{2\pi\i}(1-l/N)+\frac{\xi}{4N\pi\i}\right)\right)}
  =
  1-e^{\xi(1-l/N)}
\end{equation*}
since
\begin{equation*}
  \Re\left(\frac{\xi}{2\pi\i}\left(1-\frac{l}{N}\right)e^{\theta\i}\right)
  =
  \frac{\xi}{2\pi}\left(1-\frac{l}{N}\right)\sin\theta
  <
  \frac{1}{2}\left(1-\frac{l}{N}\right)\cos\theta,
\end{equation*}
which is between $0$ and $\cos\theta$.
Therefore we have
\begin{equation*}
\begin{split}
  \prod_{l=1}^{k}\bigl(1-e^{(N-l)\xi/N}\bigr)
  &=
  \prod_{l=1}^{k}
  \frac{\exp\left(T_{N}\left(\frac{\xi}{2\pi\i}\left(1-\frac{2l+1}{2N}\right)\right)\right)}
       {\exp\left(T_{N}\left(\frac{\xi}{2\pi\i}\left(1-\frac{2l-1}{2N}\right)\right)\right)}
  \\
  &=
  \frac{\exp\left(T_{N}\left(\frac{\xi}{2\pi\i}(1-\frac{2k+1}{2N})\right)\right)}
       {\exp\left(T_{N}\left(\frac{\xi}{2\pi\i}(1-\frac{1}{2N})\right)\right)}.
\end{split}
\end{equation*}
\par
Similarly, putting $z:=\frac{\xi}{2\pi\i}(1+l/N)$ ($0<l<N$) in \eqref{eq:lem}, we have
\begin{equation*}
  \frac{\exp
        \left(
          T_{N}
          \left(
            \frac{\xi}{2\pi\i}(1+l/N)-\frac{\xi}{4N\pi\i}
          \right)
        \right)}
       {\exp
        \left(
          T_{N}
          \left(
            \frac{\xi}{2\pi\i}(1+l/N)+\frac{\xi}{4N\pi\i}
          \right)
        \right)}
  =
  1-e^{\xi(1+l/N)}
\end{equation*}
since
\begin{equation*}
  \Re\left(\frac{\xi}{2\pi\i}\left(1+\frac{l}{N}\right)e^{\theta\i}\right)
  =
  \frac{\xi}{2\pi}\left(1+\frac{l}{N}\right)\sin\theta
  <
  \frac{1}{2}\left(1+\frac{l}{N}\right)\cos\theta,
\end{equation*}
which is between $0$ and $\cos\theta$.
So we have
\begin{equation*}
\begin{split}
  \prod_{l=1}^{k}\bigl(1-e^{(N+l)\xi/N}\bigr)
  &=
  \prod_{l=1}^{k}
  \frac{\exp\left(T_{N}\left(\frac{\xi}{2\pi\i}\left(1+\frac{2l-1}{2N}\right)\right)\right)}
       {\exp\left(T_{N}\left(\frac{\xi}{2\pi\i}\left(1+\frac{2l+1}{2N}\right)\right)\right)}
  \\
  &=
  \frac{\exp\left(T_{N}\left(\frac{\xi}{2\pi\i}(1+\frac{1}{2N})\right)\right)}
       {\exp\left(T_{N}\left(\frac{\xi}{2\pi\i}(1+\frac{2k+1}{2N})\right)\right)}.
\end{split}
\end{equation*}
Therefore we have from \eqref{eq:J_N_fig8}
\begin{equation*}
\begin{split}
  &J_N\bigl(E;\exp(\xi/N)\bigr)
  \\
  &=
  \frac{\exp\left(T_{N}\left(\frac{\xi}{2\pi\i}(1+\frac{1}{2N})\right)\right)}
       {\exp\left(T_{N}\left(\frac{\xi}{2\pi\i}(1-\frac{1}{2N})\right)\right)}
  \\
  &\quad\times
  \sum_{k=0}^{N-1}
  e^{-k\xi}
  \frac{\exp\left(T_{N}\left(\frac{\xi}{2\pi\i}(1-\frac{2k+1}{2N})\right)\right)}
       {\exp\left(T_{N}\left(\frac{\xi}{2\pi\i}(1+\frac{2k+1}{2N})\right)\right)}.
\end{split}
\end{equation*}
We use the following lemma, a proof of which is given in Section~\ref{sec:proofs_lemmas}.
\begin{lem}\label{lem:T_special}
We have
\begin{equation*}
  \frac{\exp\left(T_{N}\left(\frac{\xi}{2\pi\i}(1+\frac{1}{2N})\right)\right)}
       {\exp\left(T_{N}\left(\frac{\xi}{2\pi\i}(1-\frac{1}{2N})\right)\right)}
  =
  \frac{1}{1-e^{\xi}}
  =
  \frac{-e^{-\xi/2}}{2\sinh(\xi/2)}.
\end{equation*}
\end{lem}
\par
Now we define
\begin{equation*}
  f_{N}(z)
  :=
  \frac{1}{N}
  \left(
    T_{N}\left(\frac{\xi(1-z)}{2\pi\i}\right)
    -
    T_{N}\left(\frac{\xi(1+z)}{2\pi\i}\right)
  \right)
  -\xi z
  +2\pi\i z
\end{equation*}
so that
\begin{equation}\label{eq:Jones_sum}
  J_N\bigl(E;\exp(\xi/N)\bigr)
  =
  \frac{1}{2\sinh(\xi/2)}
  \sum_{k=0}^{N-1}\exp\left(N\times f_{N}\left(\frac{2k+1}{2N}\right)\right).
\end{equation}
Since $T_{N}(z)$ is defined for $z$ with $0<\Re(z e^{\theta\i})<\cos\theta$, the function $f_{N}(z)$ is defined for $z$ with
\begin{equation*}
  \left|\frac{\Im{z}}{\tan\theta}+\Re{z}\right|<1
\end{equation*}
from the assumption $\tan\theta<\frac{\pi}{\xi}$.
\par
From Proposition~\ref{prop:T_N_L_2}, $f_N(z)$ converges to the following function:
\begin{equation*}
  F(z)
  :=
  \frac{1}{\xi}
  \left(
    \L_2\left(\frac{\xi(1-z)}{2\pi\i}\right)
    -
    \L_2\left(\frac{\xi(1+z)}{2\pi\i}\right)
  \right)
  -\xi z+2\pi\i z
\end{equation*}
in the region
\begin{equation*}
  \left\{
    z\in\C\Bigm|
    \left|\frac{\Im{z}}{\tan\theta}+\Re{z}\right|
    \le
    1-\frac{2\pi\nu}{\xi\sin\theta},
    |\Re{z}|\le\frac{2\pi M}{\xi}-1
  \right\}.
\end{equation*}
Since $\Im\left(\frac{\xi(1\pm z)}{2\pi\i}\right)=\frac{-\xi}{2\pi}(1\pm\Re{z})$,
\begin{equation*}
  F(z)
  =
  \frac{1}{\xi}
  \left(
    \Li_2\left(e^{-\xi(1+z)}\right)
    -
    \Li_2\left(e^{-\xi(1-z)}\right)
  \right)
  +\xi z
\end{equation*}
when $-1<\Re{z}<1$ from \eqref{eq:L2}.
\begin{rem}\label{rem:F_real}
When $x$ is real and $-1<x<1$, then $F(x)$ is also real.
\end{rem}
From Lemma~\ref{lem:der_L}, the first and the second derivatives of $F(z)$ are given as follows:
\begin{align*}
  F'(z)
  &=
  \L_1\left(\frac{\xi(1-z)}{2\pi\i}\right)
  +
  \L_1\left(\frac{\xi(1+z)}{2\pi\i}\right)
  -\xi+2\pi\i,
  \\
  F''(z)
  &=
  \frac{\xi}{2\pi\i}
  \left(
    \L_0\left(\frac{\xi(1-z)}{2\pi\i}\right)
    -
    \L_0\left(\frac{\xi(1+z)}{2\pi\i}\right)
  \right)
  \\
  &=
  \frac{\xi(e^{-\xi z}-e^{\xi z})}{e^{\xi}+e^{-\xi}-e^{\xi z}-e^{-\xi z}}.
\end{align*}
Note that when $-1<\Re{z}<1$, we have
\begin{equation}\label{eq:F'}
\begin{split}
  F'(z)
  &=
  \log(1-e^{-\xi(1-z)})+\log(1-e^{-\xi(1+z)})+\xi
  \\
  &=
  \log(e^{\xi}+e^{-\xi}-e^{\xi z}-e^{-\xi z})
\end{split}
\end{equation}
from  \eqref{eq:L1}, since $\Im\left(\frac{\xi(1\pm z)}{2\pi\i}\right)=-\frac{\xi}{2\pi}(1\pm\Re{z})$.
\par
We assume that $\xi>\kappa$, where $\kappa:=\arccosh(3/2)=\log(\frac{3+\sqrt{5}}{2})=0.962\dots$.
Put
\begin{equation*}
  \varphi(\xi)
  :=
  \arccosh\left(\cosh(\xi)-\frac{1}{2}\right)
  =
  \log
  \left(
    \cosh(\xi)-\frac{1}{2}+\frac{1}{2}\sqrt{(2\cosh(\xi)-1)^2-4}
  \right).
\end{equation*}
Note that since $\xi>\kappa$, we have $\cosh(\xi)>\cosh(\varphi(\xi))$ and $\cosh(\xi)>\cosh(\kappa)=\frac{3}{2}$.
So we have $0<\varphi(\xi)<\xi$.
Note also that this definition is the same as that in Section~\ref{sec:introduction}.
\par
Since $0<\varphi(\xi)<\xi$, we have
\begin{align}
  F(\varphi(\xi)/\xi)
  &=
  \frac{1}{\xi}
  \left(
    \Li_2(e^{-\xi-\varphi(\xi)})-\Li_2(e^{-\xi+\varphi(\xi)})
  \right)
  +\varphi(\xi),
  \label{eq:F_phi}
  \\
  F'(\varphi(\xi)/\xi)
  &=0,
  \label{eq:F'_phi}
  \\
  F''(\varphi(\xi)/\xi)
  &=
  -\xi\sqrt{(2\cosh(\xi)-1)^2-4}.
  \label{eq:F''_phi}
\end{align}
The second equality follows since
\begin{equation*}
  e^{\varphi(\xi)}+e^{-\varphi(\xi)}
  =
  e^{\xi}+e^{-\xi}-1.
\end{equation*}
\par
Put
\begin{equation*}
  S(\xi)
  :=
  \xi F(\varphi(\xi)/\xi)
  =
  \Li_2(e^{-\xi-\varphi(\xi)})-\Li_2(e^{-\xi+\varphi(\xi)})
  +
  \xi\varphi(\xi).
\end{equation*}
Since
\begin{equation*}
\begin{split}
  \frac{d\,S(\xi)}{d\,\xi}
  &=
  (1+\varphi'(\xi))\log(1-e^{-\xi-\varphi(\xi)})
  -
  (1-\varphi'(\xi))\log(1-e^{-\xi+\varphi(\xi)})
  \\
  &\quad
  +\varphi(\xi)+\xi\varphi'(\xi)
  \\
  &=
  \log\frac{e^{\varphi(\xi)}-e^{-\xi}}{1-e^{-\xi+\varphi(\xi)}}
  +
  \varphi'(\xi)\log\left(e^{\xi}+e^{-\xi}-e^{\varphi(\xi)}-e^{-\varphi(\xi)}\right)
  \\
  &=
  \log\frac{e^{\varphi(\xi)}-e^{-\xi}}{1-e^{-\xi+\varphi(\xi)}}
  >0
\end{split}
\end{equation*}
and $S(\kappa)=0$, $S(\xi)>0$ for $\xi>\kappa$.
\par
Recalling that $F(x)\in\R$ when $x$ is real and $-1<x<1$, we have the following lemma.
\begin{lem}
As a function of a real variable $x$, the real-valued function $F(x)$ takes its maximum $S(\xi)/\xi$ at $x=\varphi(\xi)/\xi$ in the half open interval $[0,1)$.
Moreover, it is strictly increasing {\rm(}decreasing, respectively{\rm)} when $x<\varphi(\xi)/\xi$ {\rm(}$x>\varphi(\xi)/\xi$, respectively{\rm)}.
\end{lem}
\begin{proof}
Since we have
\begin{equation*}
  F'(x)
  =
  \log(e^{\xi}+e^{-\xi}-e^{\xi x}-e^{-\xi x}),
\end{equation*}
we see that $F'(x)$ is strictly decreasing.
Since $F(0)=0$ and $F'(x)$ becomes $0$ when $x=\varphi(\xi)/\xi<1$, the lemma follows.
\end{proof}
We have the following corollary.
\begin{cor}
For any $c_{-}$ and $c_{+}$ with $0<c_{-}<\varphi(\xi)/\xi<c_{+}<1$, there exists $\delta>0$ such that $F(x)<S(\xi)/\xi-\delta$ if $x<c_{-}$ or $x>c_{+}$.
\end{cor}
So we have $e^{N\times F(x)}<e^{N(S(\xi)/\xi-\delta)}$ if $x\not\in[c_{-},c_{+}]$.
Since $f_N(x)$ uniformly converges to $F(x)$ from Proposition~\ref{prop:T_N_L_2}, we also have $\left|e^{N\times f_N(x)}\right|<e^{N(S(\xi)/\xi-\delta')}$ for some $\delta'>0$ if $x\not\in[c_{-},c_{+}]$ and $N$ is sufficiently large.
\par
Therefore we have
\begin{align*}
  \left|\sum_{0<k/N<c_{-}}\exp\left(N\times f_N\left(\frac{2k+1}{2N}\right)\right)\right|
  &<
  c_{-}Ne^{N(S(\xi)/\xi-\delta')}
  \\
  \intertext{and}
  \left|\sum_{c_{+}<k/N<1}\exp\left(N\times f_N\left(\frac{2k+1}{2N}\right)\right)\right|
  &<
  (1-c_{+})Ne^{N(S(\xi)/\xi-\delta')}.
\end{align*}
As a result, we have
\begin{equation*}
\begin{split}
  &\sum_{k=0}^{N-1}\exp\left(N\times f_N\left(\frac{2k+1}{2N}\right)\right)
  \\
  =&
  \sum_{c_{-}\le k/N\le c_{+}}\exp\left(N\times f_N\left(\frac{2k+1}{2N}\right)\right)
  +
  O\left(Ne^{N(S(\xi)/\xi-\delta')}\right).
\end{split}
\end{equation*}
From \eqref{eq:Jones_sum}, we also have
\begin{multline}\label{eq:J_summation}
  J_N\bigl(E;\exp(\xi/N)\bigr)
  \\
  =
  \frac{1}{2\sinh(\xi/2)}
  \sum_{c_{-}\le k/N\le c_{+}}\exp\left(N\times f_N\left(\frac{2k+1}{2N}\right)\right)
  +
  O\left(Ne^{N(S(\xi)/\xi-\delta')}\right).
\end{multline}